\chapter{Environmental Considerations}

\noindent In order to be as useful as possible, FPAP describes a standard for
a wide range of systems, ideally in such a manner that FPAP could be
implemented for any programmable device with a pixel display. This, however,
comes with some considerations.

\section{Numeric Limits}

\noindent If provided an implementation shall ensure that objects of the type
specified (if provided) may represent at least the following limits specified
in the form: `\tt<type>: (<minimum>, <maximum>)\rm'.  The phrase:
`\tt maxof(<type>)\rm' refers to the maximum limmit of the type enclosed in
parenthesis.

\limitcount0\relax
\limit FPAP\_BYTE: (0, 127)
\limit FPAP\_WORD: (0, 65535)
\limit FPAP\_DWORD: (0, 4294967295)
\limit FPAP\_COLOR: (0, 1)
\limit FPAP\_BOOL:  (FPAP\_FALSE, FPAP\_TRUE)

\noindent The type `FPAP\_PTR' has unspecified limits, yet shall be able to
store the addresses of all objects used as arguments to fpap functions and
those of all objects generated by fpap functions.

The type `FPAP\_ERROR' shall be able to represent the maximum number of error
codes referenced by a single FPAP function (\S TODO) (these codes are listed
in \S 2.2 Error Codes).

The type `FPAP\_PROPERTIES' has limits implementation defined according to the
number of properties defined for the `fpap\_set' and `fpap\_get' functions.

As such the following should be defined within the documentation for the
implemtnation, or otherwise specified within the program with the identifiers
given (i.e., C macros):

\halign{\tt#\hfil\cr
	FPAP\_BYTE\_MAX >= 127,\cr
	FPAP\_WORD\_MAX >= 65535,\cr
	FPAP\_DWORD\_MAX >= 4294967295,\cr
	FPAP\_COLOR\_MAX >= 1,\cr
	FPAP\_BOOL\_MIN == FPAP\_FALSE,\cr
	FPAP\_BOOL\_MAX == FPAP\_TRUE,\cr
	FPAP\_PTR\_MIN == implementation defined,\cr
	FPAP\_PTR\_MAX == implementation defined,\cr
	FPAP\_ERROR\_MIN == defined by error codes,\cr
	FPAP\_ERROR\_MAX == defined by error codes,\cr
	FPAP\_PROPERTIES\_MIN == implementation defined,\cr
	FPAP\_PROPERTIES\_MAX == implementation defined\cr
}
	

\section{Error Codes}

\noindent The following identifiers shall specify values according to errors
referenced in FPAP functions. The specific numeric values are implementation
defined.

\error FPAP\_SUCCESS: Defines the success condition for FPAP functions, all
objects of the \tt FPAP\_ERROR \rm type must be able to represent
\tt FPAP\_SUCCESS \rm and compare equal with any other object of that type with
the value \tt FPAP\_SUCCESS\rm.
\error FPAP\_NO\_WINDOW\_SYSTEM: An error occurred while trying to access a
function of the window system.
\error FPAP\_NO\_WINDOW: An error occurred while trying to access a function
of the region specified by the framebuffer.
\error FPAP\_WINDOW\_CLOSED: The window has been closed and it is undefined
behaviour to access the framebuffer.
\error FPAP\_UNDEFINED\_PROPERTY: The property specified does not exist.
\error FPAP\_INACCESSIBLE: The property specified could not be accessed.

\section{Properties}

\noindent A property is an implementation defined value used to index objects
within the fpap instance. They may be get and set by the user with the
respective \tt fpap\_get \rm and \tt fpap set \rm commands.

Identifiers used to designate properties shall be prefixed with
`FPAP\_PREFIX\_'. Such that if an implementation should have a property for
vsync, it may be called `FPAP\_PREFIX\_VSYNC'.

\vfill\break
